% @file report.tex
% @brief Пояснительная записка курсового проекта №18 «MinIO Object Storage».
% @author Тимошенко (заполнить ФИО)
% @date 2025-05-25
% @version 1.0.0

\documentclass[14pt,a4paper]{extreport}
\usepackage[utf8]{inputenc}
\usepackage[T2A]{fontenc}
\usepackage[russian]{babel}
\usepackage{indentfirst}
\usepackage{graphicx}
\usepackage{hyperref}
\usepackage{listings}
\usepackage{geometry}

\geometry{
  a4paper,
  left=30mm, right=15mm,
  top=20mm, bottom=20mm
}

\title{Пояснительная записка\\к курсовому проекту №18\\
\large Развёртывание и настройка распределённого хранилища объектов MinIO}
\author{Тимошенко (заполнить ФИО)}
\date{2025}

\begin{document}

\maketitle
\tableofcontents

\chapter{Введение}

Актуальность: рост объёмов неструктурированных данных требует масштабируемых объектных
хранилищ с поддержкой S3-совместимого API. MinIO — лидирующее open-source решение класса
Self-Hosted Object Storage.

Цель работы: развернуть и настроить производственный кластер MinIO с обеспечением
отказоустойчивости (Erasure Code), безопасности (TLS, Docker Secrets) и наблюдаемости
(Prometheus + Grafana).

\chapter{Анализ существующих решений}

\section{Сравнительный анализ}

Для реализации объектного хранилища рассматривались два подхода:
\begin{enumerate}
  \item \textbf{MinIO} — S3-совместимое хранилище с нативной поддержкой Erasure Code,
        развёртывание через Docker, встроенные метрики Prometheus.
  \item \textbf{Ceph RADOS Gateway} — распределённая СХД, поддерживает S3 API,
        но требует значительно больших ресурсов и экспертизы для настройки.
\end{enumerate}

\textbf{Выбор: MinIO} обоснован минимальными требованиями к ресурсам,
полной совместимостью с AWS S3 SDK и простой интеграцией с Docker Compose.

\chapter{Архитектура решения}

Кластер состоит из 4 нод MinIO в режиме Erasure Code (8 дисков = 4$\times$2).
Трафик терминируется на Nginx с TLS. Мониторинг — Prometheus + Grafana.

\chapter{Описание компонентов}

\section{MinIO Erasure Code Cluster}
Четыре контейнера minio1--minio4 образуют единый кластер с алгоритмом Reed-Solomon.
Кластер переживает одновременный выход из строя до 4 дисков (n/2).

\section{Nginx TLS Proxy}
Nginx выполняет TLS-терминацию и балансировку нагрузки методом \texttt{least\_conn}.

\section{Мониторинг}
Prometheus собирает метрики с эндпоинта \texttt{/minio/v2/metrics/cluster}.
Grafana визуализирует дашборды, хранящиеся в Git как JSON.

\chapter{Процесс развёртывания}

Развёртывание одной командой:
\begin{verbatim}
sudo bash deploy/scripts/bootstrap.sh
cp .env.example .env && nano .env
cd deploy && docker compose up -d --build
\end{verbatim}

\chapter{Результаты тестирования}

Тестирование включает:
\begin{itemize}
  \item Проверку healthcheck всех контейнеров (\texttt{test\_healthcheck.sh})
  \item Интеграционные тесты доступности API (\texttt{test\_minio.sh})
  \item HA-тест: отключение одной ноды при сохранении доступности кластера
\end{itemize}

\chapter{Заключение}

Разработана и развёрнута полноценная инфраструктура объектного хранилища MinIO.
Достигнуты все цели: отказоустойчивость (Erasure Code), безопасность (TLS, Docker Secrets,
hardening ядра), наблюдаемость (Prometheus, Grafana). Вся инфраструктура разворачивается
из репозитория одной командой.

\end{document}
